\section{The fixed port and its robust decision gadget}\label{sec:inner-gadget}

We now prove Theorem~\ref{thm:port-gadget}. Throughout this section all
fields, groups, block lengths, and path labels are those of the inner
construction of Sections~\ref{sec:inner-algebra} and~\ref{sec:inner-binary}.
In particular, its polynomial-size affine group is unrelated to the outer
index group. We first construct the port independently of the eventual
arity. We then fix an actual arity and construct a generalized subcircuit
for one Boolean decision. The subcircuit will use the geometric procedures
of Lemma~\ref{lem:geometric-procedures}, whose proof is independent of this
section. Word distances in this section are relative Hamming distances;
vector distances use the displayed normalized block norm.

\subsection{A port fixed before the arity}

Fix the exponent allowance $c_{\rm allow}$ and the common field, point,
slot, and row-code parameters supplied by
Lemmas~\ref{lem:inner-algebra-interface} and
\ref{lem:inner-binary-interface}. In particular these parameters are
chosen from $n'=(L+2)^{c_0}$, with fixed $c_0>c_{\rm allow}+5$, before
specifying the number of input ports or the actual Boolean circuit.
Construct the one-port instance of the inner path with the prescribed
initial copy phase and, afterwards, a fixed trivial computation. Include
all its positive-tag validity, critical-address, derivative, and
verification-auxiliary families, as well as its tag-zero families. This is
the \emph{dummy template}. Write $a_s$ for its distinguished source and
$a_s^+$ for its immediate successor. Their common tape has blank work,
stage-zero legacy, input $s$, frame $\gamma^{-1}$, and post mode. Their
progress values are respectively $q-1$ and $q-2$. Define
\begin{equation}\label{ig:port-definition}
 \mathcal D(s)=\bigl(\mathcal W_{\rm dum}(a_s),
                         \mathcal W_{\rm dum}(a_s^+)\bigr),
 \qquad S_0=2m_{\rm dum}.
\end{equation}
Both halves use the entire balanced binary word of the same dummy template.

\begin{proposition}[Fixed dummy port]\label{prop:dummy-port}
The map in~\eqref{ig:port-definition} has polynomial length, constant
relative Hamming distance $d_{\mathcal D}>0$, and the following properties.
\begin{enumerate}[label=\textup{(\roman*)}]
\item Every output bit is a multilinear polynomial in $s$ over $\F_2$ of
 degree at most six. Its coefficients, and the whole canonical encoding,
 are computable deterministically in time polynomial in $L$.
\item There is a perfectly complete tester using a bounded number of bit
 reads and bounded Boolean processing. Its public randomness is uniformly
 enumerable with polynomial-size support and rational probabilities of
 polynomial bit length. Every fixed branch has unconditioned read
 subdistributions of bounded density relative to $[S_0]$. There is a
 constant $C_B$ such that proximity $a$ to a port word implies rejection
 at most $C_Ba$. For every fixed $a>0$, there is a fixed
 $\theta_B(a)>0$ such that rejection less than $\theta_B(a)$ implies
 proximity less than $a$ to some $\mathcal D(s)$, for all sufficiently
 large $L$.
\item There are a deterministic polynomial-time exact decoder at a fixed
 positive Hamming radius and a fixed
 $0<\nu_B<d_{\mathcal D}/4$. From any word within $\nu_B$ of
 $\mathcal D(s)$, either half supplies bits of its canonical word with
 any fixed requested average accuracy under any fixed bounded-density
 law on the original edge and outer indices. For a product block this
 law is on the joint pair of indices. Requested masks and additive
 forms may depend arbitrarily on the public request. The read procedures
 have unconditioned fixed-branch smoothness, including their unsuccessful
 branches.
\end{enumerate}
The dummy template, its decoder radius, and $\nu_B$ are all fixed before
any later actual arity, circuit-size factor, or correction-round count.
They are not decreased when those constants increase.
\end{proposition}

\begin{proof}
The inner distance assertion applies separately to each half. Distinct
inputs give distinct labels already in the order-one input track, so it
gives the claimed distance for their concatenation.

For testing, mix the ordinary valid-label testers on the two halves with
the following additional checks. At systematic point probes compare the
legacy chart with the public stage-zero polynomial, the mode with post,
and the progress marker with its specified one-hot tensor, using the
source or successor progress value as appropriate. Uniform slot masks
and additive binary forms give smooth bit versions of these checks; a
specified form can also be split into fresh uniform summands. On exact
valid labels, an incorrect location condition has a fixed positive
rejection gap by the low-degree or constant-tensor distance. Correct
conditions force the blank-work tape, the specified progress and mode,
and frame $\gamma^{-1}$. Finally compare the two order-one input point
tracks at identical original systematic edge/point indices and identical
uniform bit forms, choosing copies independently in the two halves.
Different inputs give a fixed positive discrepancy gap. The two reads
have smooth marginals individually; they do not require a product block.
Each mixture weight is a fixed positive constant.

Here is the fixed-gap implication explicitly. Let $g>0$ be a common lower
bound for the rejection of any violated location or equality requirement
on exact pairs of valid labels, with the mixture weights included. Let
$C$ bound the change in these tests' rejection under changes to a fraction
of the whole port. Given $a>0$, choose
\[
 0<a_0<\min\{a,g/(4C)\}.
\]
Ordinary soundness supplies a positive rejection threshold forcing each
half within $a_0$ of an exact valid label. Take the total tester threshold
small enough to imply both of these assertions and also smaller than
$g/2$. The pair of nearby exact labels is at distance less than $a_0$
from the received port. If any required exact condition were wrong,
smoothness would make the received rejection at least
$g-Ca_0>g/2$, a contradiction. Thus the exact pair is precisely
$\mathcal D(s)$ and the received distance is less than $a$. This defines
$\theta_B(a)$. Perfect completeness and the upper bound $C_Ba$ follow by
charging every read whose answer differs from the exact port word. All
choices and computations have the bounded, polynomially enumerable
properties asserted in the two inner interfaces.

At the two fixed labels all logical inputs except $s$ are public.
Canonical persistent values and their coordinate coefficients have
logical degree at most three, even on invalid Boolean tapes. Their
verification certificates use public linear maps on one primary array:
multiplication by the public test matrix, summation, division by a public
monic polynomial, coefficient extraction, and interpolation. These
operations preserve degree. Fixed-progress stack formation, restriction
to a point or a line, line-coefficient extraction, slot contraction, row
extension, and additive bit forms are linear over the prime field. Thus a
$Y$ bit has degree at most three. A $Z$ bit is a binary bilinear form of
two such operands and has degree at most six. Repeating bits for balance
changes nothing. The degree bound is independent of the allowance and
of all future template choices.

For completeness, the coefficient algorithm need not manipulate a
large symbolic expression. For a specified bit function $b$, evaluate
the canonical dummy word on $1_T$ for all $T\subseteq[L]$ with
$|T|\le6$. Its coefficient at the squarefree monomial indexed by $A$ is
\[
 c_A=\sum_{T\subseteq A}b(1_T)\quad\text{in }\F_2,
 \qquad |A|\le6.
\]
Boolean multilinear inversion and the degree bound give every
coefficient; there are $O((L+1)^6)$ evaluations, each polynomial-time.

Use the deterministic order-one recovery from
Lemma~\ref{lem:inner-binary-interface} on the first half. A port error
fraction $r$ causes an error fraction at most $2r$ there, so a fixed
sufficiently small $r$ gives exact recovery of $s$. Finally choose
$\nu_B<d_{\mathcal D}/4$ with $2\nu_B$ inside the dummy average-correction
radius and, if necessary, the just established recovery radius.
A port within $\nu_B$ corrupts either half by at most $2\nu_B$.
The dummy corrector therefore answers the stated requests from that
half; its read density relative to the whole port loses only a factor
two. Its radius is independent of the requested accuracy and density
constant. Every constant used here belongs to the one-port template.
\end{proof}

\subsection{Exact original-index correspondence and virtual source bits}

Now fix an actual arity $k$, a circuit-size factor, and a Boolean circuit
$f$ covered by Theorem~\ref{thm:port-gadget}, padding unused inputs with
public zero ports. Fix words $y_i$ with
\begin{equation}\label{ig:fixed-ports}
 \dist\bigl(y_i,\mathcal D(s_i)\bigr)\le\nu_B.
\end{equation}
For a public zero port all bits of $\mathcal D(0^L)$ are computed publicly.
The true graph used below is the single actual-template path for the
fixed tuple $s=(s_1,\ldots,s_k)$, rather than the union over possible
input tuples. Write $w_*(s)$ and $w_*^+(s)$ for its source and immediate
successor. These have the same public frame, tape locations, mode, and
progress as their dummy counterparts, but their input-dependent encoding
bits must be supplied from the ports.

\begin{lemma}[Source correspondence and virtual access]\label{ig:virtual-source}
For either $w_*(s)$ or $w_*^+(s)$, every tag-$i$ point or line $Y$
operand agrees with its dummy counterpart for $s_i$ at the same original
indices, after renaming the port tag. The same is true of a $Z$ block
whose two operands both have tag $i$. Tag-zero data at either label are
publicly computable. Consequently source and successor encoded-bit
requests with bounded original-index density admit arbitrarily accurate
average answers from~\eqref{ig:fixed-ports}, with smooth external reads.
For mixed products the requisite density is joint density on the complete
edge/point indices before marginalization. Actual and dummy balanced-copy
multiplicities need not agree.
\end{lemma}

\begin{proof}
Apply Lemma~\ref{ia:source-data} at the two labels in
\eqref{ia:source-successor}. Both have frame $\gamma^{-1}$, blank work,
stage-zero legacy, and post mode; their progresses are respectively
$q-1$ and $q-2$. At either progress the lemma identifies each positive
algebraic array with its one-port counterpart at the same native point,
slot, and mask indices. This includes verification auxiliaries and all
derivative descendants: Lemma~\ref{ia:port-ancestry} fixes their separate
call layouts independently of work destinations, the number of ports,
and the actual circuit. The placements and fixed-progress stack formulas
are common as well.

The common row field, row matrices, systematic extension, and point/line
domains were fixed before $k$. Applying these same linear operations
preserves the equality at the original encoded indices, for the source
and its first decrement separately. The actual template may balance its
block types with different copy counts. Forgetting its copy index and
choosing a fresh uniformly random dummy copy preserves the original index
and costs only the fixed actual block-weight factor.

The same source-data lemma makes every tag-zero array public at these
labels. Its conclusion includes derivatives evaluated on public shifted
work, even when the shifted Boolean tape is invalid: canonical certificate
recipes are defined there and do not solve a final validity equation.
Thus no work-side certificate needs to evaluate $f(s)$ or a stored sum
with a port-dependent primary. The actual encoded $J$-symbols of these
public arrays, including their parity values, are computed with the
actual template's common row encoder.

The only remaining case is a product with tags $(i,0)$ or $(0,i)$.
Write $x$ for the positive operand's entire original edge/point index and
$z$ for the public operand's entire index. After the indicated contractions
and public scalar multipliers its bit is, in one orientation,
\begin{equation}\label{ig:mixed-form}
 Q\bigl(\pi(b_i u(x)),\pi(b_0 v_{\rm pub}(z))\bigr)
       =\mu_{z,b_i,b_0,Q}\bigl(u(x)\bigr),
\end{equation}
where $\mu_{z,b_i,b_0,Q}$ is a specified additive binary form. The dummy
$Y$ corrector permits precisely such arbitrary forms, including dependence
on both requested positions and all the public coefficients. It retains
$x$ at the same original index. If the request law satisfies
\[
 p(x,z)\le\frac{C}{|\mathcal X|\,|\mathcal Z|},
 \qquad\text{then}\qquad
 \sum_zp(x,z)\le\frac{C}{|\mathcal X|}.
\]
This proves the required density for the dummy $Y$ request on both its
edge and its point. A same-positive-tag $Z$ request retains both original
indices and their joint density; even when its types coincide, no pair
is replaced by a diagonal. A $(0,0)$ product is public. There are no
products of two distinct positive tags. Proposition~\ref{prop:dummy-port}
now supplies all the answers. Fixing a source/successor choice, orientation,
or block type costs a fixed factor. All error and density bounds are
summed before any later answer-dependent selection; no density bound is
conditioned on that selection.
\end{proof}

\subsection{Validation of the true path and numerical access}

Let $m=m_{\rm in}$ be the length of one actual-template word, and let
$b_\#$ be its number of equally weighted original binary block types.
The path labels have fixed separation $d_*>0$ in the normalized Euclidean
norm. Constants from this point may depend on the actual arity,
circuit-size factor, and desired output-loss fraction $\upsilon>0$.
They will not alter $\mathcal D$ or $\nu_B$.

The ordinary binary tester recognizes the union of paths over all input
tuples. Before using a seed for geometry, we must ensure that its nearby
label belongs to the path for the externally supplied tuple $s$.
We do this by comparing each input track with its corresponding port.

\begin{lemma}[Validation at the fixed external tuple]\label{ig:validation}
There are fixed $\rho_{\rm in},\rho_{\rm in,comp}>0$ and bounded-read
validation routines for an internal binary seed $z\in\{0,1\}^m$ such
that, with arbitrarily small prescribed failure probability, acceptance
implies distance less than $\rho_{\rm in}$ to a unique label on the
true-$s$ path, and distance less than $\rho_{\rm in,comp}$ to a true-path
label implies acceptance. These radii are chosen before the reliability
of the routines. For each fixed seed that can be accepted on good
decisions its reference label is fixed independently of the decision
randomness. Current and actual incident-neighbor bits and existence flags
can subsequently be supplied with arbitrary fixed average accuracy for
the request laws of Lemma~\ref{lem:geometric-procedures}.
\end{lemma}

\begin{proof}
Choose $\rho_{\rm in}$ within the whole-word uniqueness and actual-template
correction radii. Estimate rejection of the ordinary tester for the union
of all port tuples. For each port $i$, also sample the order-one point
track $W_{s_i}$ of $z$ at uniform systematic original edge/point indices
and uniform forms, and compare it with the bit at the identical original
index of the corrected dummy source half of $y_i$. Choose a random actual
copy for the direct read. This input track is unrotating and uses the
same extension at every actual path label.

On an exact valid label the ideal mismatch is zero if its $i$th port is
$s_i$, and at least a constant $p_s>0$ otherwise: the nonzero order-one
polynomial differs on a constant fraction of represented arguments, and
a uniform additive form detects a nonzero difference with probability
$1/2$. We can choose one $p_s$ for the fixed templates. A balanced error
fraction $a'$ in $z$ changes a comparison rate by at most $C'_k a'$,
where $C'_k$ includes the systematic and block-weight density factors.
Choose
\[
 0<a'<\rho_{\rm in},\qquad C'_k a'<p_s/8.
\]
Use half the ordinary soundness threshold for $a'$ as the rejection
cutoff, and $p_s/2$ as each mismatch cutoff. Good statistics with fixed
positive slack imply that an accepted seed is $a'$-close to an exact
valid label. The mismatch gap then forces every port of that label to
be $s_i$. Conversely, smoothness of the ordinary tester and of the
comparison probes makes all rates lie strictly below their cutoffs when
the seed is within a sufficiently small fixed
$\rho_{\rm in,comp}>0$ of a true-path label.

First fix sample counts that concentrate the underlying ideal rates at
the required slack; independent samples and Chebyshev's inequality
suffice. Next choose the dummy correction accuracy so that the bounded
number of supplied comparison bits fail with the desired total
probability. This order leaves $a'$, $\rho_{\rm in,comp}$, and $\nu_B$
unchanged. On good acceptance, the ordinary correction and neighbor
procedures of Lemma~\ref{lem:inner-binary-interface} apply to the fixed
unique reference. They return the true incident flags, including the
disabled wrap flag, and true encoding bits at existing neighbors with
arbitrary average accuracy. Analyze each possible reference of the fixed
seed before intersecting with acceptance; one never conditions a
smoothness estimate on that event. On unsuccessful branches every routine
still has a bounded Boolean output or a fixed default.
\end{proof}

The point-vector space for the gadget is
\begin{equation}\label{ig:point-space}
 \mathcal B=[-1,2]^{4m},\qquad
 \|U\|_{\rm b}^2=\frac1m\sum_{i=1}^{4m}U_i^2.
\end{equation}
Among its first-block coordinates select the systematic point $Y$ bits
of the constant output track $F_{\rm out}$, with forms $\mu$ satisfying
$\mu(1)=1$, including all their balanced copies. At the unique sink all
these bits equal $f(s)$. At least half the row labels are systematic,
half the forms take $1$ to $1$, and the track has weight $1/b_\#$.
Consequently the selected fraction among all $4m$ coordinates is at least
\begin{equation}\label{ig:selection}
 c_{\rm sel}=\frac1{16b_\#}.
\end{equation}
Choose the geometric parameter $\tau>0$ so small that
\begin{equation}\label{ig:tau}
 \frac{(5\tau)^2}{(2/5)^2}<\rho_{\rm in,comp},
 \qquad
 \frac{16(10\tau)^2}{c_{\rm sel}}<\frac{\upsilon}{4}.
\end{equation}
A binary disagreement after rounding a physical block at any threshold
in $[2/5,3/5]$ incurs squared error at least $(2/5)^2$ from the true bit.
Thus the first inequality gives the geometric completeness condition.
Obtain fixed $B_1,c_*,\beta_0>0$ from
Lemma~\ref{lem:geometric-procedures} and its geometric lower bound,
increasing $B_1$ to an integer if necessary. Write $\mathcal E$ for the
endpoint set of the true path. It excludes the top of the distinguished
vertical segment; here every member of $\mathcal E$ has first block equal
to the unique sink encoding. In particular,
\begin{equation}\label{ig:geometric-bound}
 \|v(U)\|_\infty\le B_1-1,
 \qquad
 \|\clip_{\mathcal B}(U+\beta v(U))-U\|_{\rm b}\ge c_*\beta
\end{equation}
whenever $\dist(U,\mathcal E)>10\tau$ and $0<\beta\le\beta_0$.

We verify the numerical access contract rather than presupposing access
to hidden circuit values. There are at most eight true word roles: each
of two seed roles and its two possible neighbors, and the special source
and successor. Four incident edge roles contribute four straight pieces
and three internal fillets each; with two possible seed junctions and
the special vertical and its bottom fillet there are at most 32 piece
roles. Lemma~\ref{ig:virtual-source} always supplies the special roles,
even when neither seed validates. The other roles use
Lemma~\ref{ig:validation} and the actual neighbor routines. Their requests
are full-word uniform samples for statistics or the coordinate matching
the requested field component. Condition on a fixed output block type,
or sum over the finite type list, to obtain the required edge/outer laws,
including joint laws for $Z$ blocks. Mode, progress, and critical decisions
use their own uniform systematic marginals.

Fix the geometric mesh, numerical tolerances, scalar sample counts, and
high-level role list first. Allocate calls for the possible roles and
matching components before any numerical minimum or branch is selected.
For a fixed input and rounding threshold, each seed has at most one
reference that good validation could accept, so there are only boundedly
many underlying actual candidates. Independent uniform-coordinate
samples concentrate all their bounded statistics, including physical
block statistics, without conditioning on which candidate wins. Set the
required bit and decision error budgets for these high-level calls next,
and only then fix the lower-level correction and repetition counts.
A numerical winner chooses a reserved role and mesh case; it cannot turn
a hidden answer into a new arbitrary word address.

Raw scalar samples can be quantized on fixed-width grids with public
random offsets. Their finite range and fixed widths give finitely many
answers. All public field, address, coefficient, mesh, reciprocal,
square-root, and trigonometric calculations are made during the
reduction, to the fixed required precision. A finite table indexed by
the bounded answers supplies the required numerical operations. The
procedures return bounded values on arbitrary unsuccessful trials as
well. Refining an offset grid later changes the thresholds and public
table entries, but not the bin count, word-address primitives, or number
of answer slots.

It follows from Lemma~\ref{lem:geometric-procedures}, with these explicit
oracles, that for $\kappa=1/100$ we can fix a finite unscaled estimator
whose signed output has the form
\begin{equation}\label{ig:ideal-estimator}
 \widehat v_i(U,y)=B_1(p_i^*-n_i^*),\qquad
 p_i^*,n_i^*\in[0,1]\text{ fixed-precision dyadics},\quad p_i^*n_i^*=0,
\end{equation}
and, whenever $\dist(U,\mathcal E)>10\tau$,
\begin{equation}\label{ig:ideal-error}
 4\E_{i,t}\bigl|\widehat v_i(U,y)-v_i(U)\bigr|^2
        <(c_*\kappa/100)^2.
\end{equation}
Here $i$ is uniform on $[4m]$, $t$ is the true public tape, and the
inequality holds conditional on each fixed admissible internal rounding
threshold and scalar offset. Rounding the signed output to its fixed
dyadic grid is included in the accuracy. All choices to this point are
independent of $\beta$ and of any external comparison margin.

\subsection{Full tapes, deterministic copies, and read routing}

We require bounds for every read actually wired into the circuit. Unfold
all hypothetical Boolean, validation, equality, and quantization answers,
assigning fresh public randomness to unexecuted nodes as well. After the
accuracy choices above this is a fixed finite tree. Let $H$ bound the
number of its raw comparison occurrences. A branch may select an
unsuccessful seed, a wrong flag, a wrong transition type, or a default;
these choices do not license a new addressing primitive.

More explicitly, list each potential initial row-cone path before
removing repeated positions. On a frozen symbolic branch the covariance
map $(g,I)\mapsto(gh,h^{-1}I)$ is a joint permutation; a further row-star
step is a fixed left multiplication of the edge coordinate, with a change
among boundedly many edge labels. On product blocks the two axes undergo
these transformations separately. Neighbor proposals use a finite list
of the affine pullbacks from the inner algebra. Omissions, deduplication,
and representative choices only select subdistributions of the listed
paths, without renormalization. A flagged zero is a public value and
needs no artificial dummy read at a concentrated address. Information-set
inverses and linear-expression coefficients are public functions of the
frozen branch and indices. Fresh additive splitting is performed after
all form or multiplier pullbacks, so arbitrary answer-dependent requested
forms still give the uniform inner marginals. These facts are exactly
the unconditioned guarantees of
Lemma~\ref{lem:inner-binary-interface}; they apply whether the answers
on the frozen branch are correct or not.

These inner routines satisfy the additional physical-read
hypothesis~\ref{geo:access-finite} of
Definition~\ref{geo:access-contract}. We therefore invoke part~(iv) of
Lemma~\ref{lem:geometric-procedures} at the physical-read level: each
completed inner bit or decision already has a bounded physical answer
tree with the stated unconditioned branch marginals.

For a fixed internal offset choice, sum these bounds over all potential
slots and geometric roles. With $i$ uniform on $[4m]$, the expected number
of reads of a specified internal base coordinate is at most $C/(4m)$,
and of a specified position of any one external port is at most $C/S_0$,
for a fixed constant $C$. Internal direct reads use uniform statistic
positions or matching output positions. Virtual external reads use the
original-index correspondence and the marginal calculation
in~\eqref{ig:mixed-form}. The constants $H,C$ are uniform over the
function and over the external threshold $\xi$ for the fixed bounds,
and remain unchanged by later offset-grid refinement. These are counts
on the whole tree, not smoothness conditioned on execution.

Choose finite uniform rational grids for the internal common binary
threshold $\xi_{\rm in}\in[2/5,3/5]$ and the scalar offsets modulo the
already fixed quantization widths. If $D_{\rm grid}$ is a lower bound
on their sizes, then for a fixed scalar value the probability of being
within $4\omega_{\rm in}$ of a binary threshold or one of its relevant
shifted boundaries is at most
\begin{equation}\label{ig:ambiguity}
 O\bigl(D_{\rm grid}^{-1}+C_{\rm width}\omega_{\rm in}\bigr),
\end{equation}
where $C_{\rm width}$ is fixed by the widths and finite range. Sum the
conditional slot-density bounds over all base positions at every offset,
and then over the bounded slot list. Thus the true-law fraction of cones
having any such ambiguous internal comparison is uniformly as small as
desired, for every fixed mean array, by refining the grids and decreasing
the positive rational $\omega_{\rm in}$. External comparisons are not
included here; their spoil budget is a separate hypothesis.

For each output coordinate $i$, the full tape law, including the offsets
and conditional random choices at all hypothetical nodes, has enumerable
polynomial support and rational probabilities of polynomial bit length.
A fixed product or finite tree of polynomial-size choices still has
these bounds. Choose a common square $R$ and replace each law by an equal
list of $R$ tapes with total variation at most
\begin{equation}\label{ig:tv}
 t_{\rm tv}=\frac{\eta_{\rm tv}}{\max\{4m,S_0\}}.
\end{equation}
Indeed floor $R$ times each probability and distribute the remaining
integer counts within the support. If $P(L)$ is a common support bound,
choosing a square $R\ge P(L)\max\{4m,S_0\}/\eta_{\rm tv}$ suffices.
A square within a fixed factor of this polynomial bound is available,
and the integer and rational operations cost polynomial time. The bound
and $R$ can be common to all padded functions in the stated class.
Use the $r$th listed tape at physical replica $(i,r)$ and put
\begin{equation}\label{ig:gscale}
 G_0=4mR.
\end{equation}
By bounded outputs, sufficiently small fixed $\eta_{\rm tv}$ preserves
from~\eqref{ig:ideal-error} the empirical bound
\begin{equation}\label{ig:empirical-error}
 \left(\frac1{mR}\sum_{i,r}
       |\widehat v_{ir}(U,y)-v_i(U)|^2\right)^{1/2}
       \le c_*\kappa/20,
\end{equation}
as well as the desired rare-event slack in~\eqref{ig:ambiguity}.
This estimate holds for every fixed $U$: the tape lists depend only on
the public laws, and total variation controls every bounded function of
the tape, including its error on that $U$.

For clarity, the unequal-size load calculation is
\begin{align}
 \#\{\text{reads of external position }j\}
 &\le G_0\left(\frac C{S_0}
            +\frac{H\eta_{\rm tv}}{\max\{4m,S_0\}}\right)
 \le (C+H\eta_{\rm tv})\frac{G_0}{S_0},
 \label{ig:external-load}\\
 \#\{\text{reads of internal base coordinate }j\}
 &\le G_0\left(\frac C{4m}
            +\frac{H\eta_{\rm tv}}{\max\{4m,S_0\}}\right)
 \le (C+H\eta_{\rm tv})R.
 \label{ig:internal-load}
\end{align}
For one position the count on a tape lies in $[0,H]$, so changing its
law by total variation $t_{\rm tv}$ changes its expectation by at most
$Ht_{\rm tv}$, which proves these inequalities. They hold even when
$m/S_0$ grows as an arbitrarily large fixed power of $L$.
Route the internal occurrences round-robin to the $R$ replicas of their
base coordinate. Every replica is then read a bounded number of times.
Every external occurrence is exposed separately for the caller to route;
\eqref{ig:external-load} proves the constant $C_{\rm occ}$ in
Theorem~\ref{thm:port-gadget}.

\subsection{An explicit generalized subcircuit}

Let $R=b^2$ and identify replica labels with $(\mathbb Z/b\mathbb Z)^2$.
On each base coordinate use the average $A$ of the eight labeled moves
\[
 (u\pm2v,v),\quad (u\pm(2v+1),v),\quad
 (u,v\pm2u),\quad (u,v\pm(2u+1))
\]
and eight identity labels, each with weight $1/16$. The all-modulus
Margulis--Gabber--Galil bound, in the primary formal development
\cite[\texttt{mgg\_numerical\_radius}]{MGGAFP}, bounds the normalized
upper mean-zero eigenvalue of the eight-move operator by $5\sqrt2/8$.
The operator is symmetric with spectrum in $[-1,1]$. Half-laziness
therefore yields
\begin{equation}\label{ig:replica-gap}
 \|AD\|\le(1-\kappa)\|D\|,\qquad \kappa=1/100,
\end{equation}
on replica deviations, since $(1+5\sqrt2/8)/2<1-\kappa$.
The statement includes every positive integer $b$ and retains loops
and repeated labels. No prime-modulus choice or randomized graph search
is needed.

Each physical replica is a node $x_{ir}\in[0,1]$, interpreted as
$X_{ir}=3x_{ir}-1\in[-1,2]$. Its own \emph{direct cone} produces it
using the prescribed 16 averaging inputs and its fixed estimator tape.
All gates of different direct cones are distinct; they may read the
same existing node as input. The estimator uses only the nine gate types
of Definition~\ref{def:gcircuit}, as follows.

For every raw comparison occurrence create its own Constant threshold
node and a Less gate, oriented to give the desired high bit. Internal
binary thresholds and quantization boundaries are translated from
physical coordinates by $a\mapsto(a+1)/3$. Out-of-range constant cases
use Boolean constants. An external occurrence reads its individually
routed external node and compares it with a Constant at $(\xi+1)/3$.
Threshold gates are inside the cone and are never shared between cones.
All public zero bits are produced by Constant gates.

Unfold finite answer-dependent choices and implement their Boolean
answers and finite lookup outputs by truth-table circuits over And, Or,
and Not. A multiplexer, for example, is
$(c\mathbin{\mathrm{And}}a)\mathbin{\mathrm{Or}}
 ((\mathrm{Not}\,c)\mathbin{\mathrm{And}}b)$.
Copy gates may provide fan-out, although unrestricted input fan-out
already suffices. Induction through these circuits preserves a correct
near-bit with the same tolerance: weak Boolean antecedents apply to
inputs at most $\eps$ or at least $1-\eps$. This costs a fixed number
of gates because the complete answer tree and output precision have
already been fixed. In particular, polynomial-size field elements and
indices are computed publicly to choose wires, coefficients, and tables;
they are not hidden strings being processed by a bounded Boolean cone.

Apply Scale with parameter $1/16$ to each averaging input and accumulate
with Add to obtain their mean. Use Scale and Add to sum the fixed dyadic
weights of the signed-output bits, obtaining nodes $p,n\in[0,1]$ for
$p^*,n^*$ in~\eqref{ig:ideal-estimator}. For a positive rational $\beta$
chosen below, put $t_\beta=\beta B_1/3\le1$. Scale $p,n$ by
$t_\beta$, add the positive term to the averaging node, and subtract
the negative term, in that order, to produce $x_{ir}$. All parameters
lie in $[0,1]$. Exact partial sums for averages and dyadic expansions
are at most one; hence their clipped Add realization is exact on ideal
inputs. Clipped sums and differences are Lipschitz in their inputs, so
there is a fixed arithmetic constant $C_{\rm ar}$ with the following
two implications whenever the entire direct cone is $\eps$-satisfied.
First,
\begin{equation}
 |X_{ir}-(AX)_{ir}|
   \le2\beta B_1+C_{\rm ar}\eps.
       \label{ig:coarse-direct}
\end{equation}
Under the same hypothesis, if all its raw comparisons give the ideal
near-bits, then
\begin{equation}
 \left|X_{ir}-\clip_{[-1,2]}
       \bigl((AX)_{ir}+\beta\widehat v_{ir}\bigr)\right|
   \le C_{\rm ar}\eps.
       \label{ig:fine-direct}
\end{equation}
Here are explicit error counts. Sixteen Scale gates and fifteen Add
gates produce an averaging node $\widetilde a$ with
$|\widetilde a-a|\le31\eps$, where $a=(Ax)_{ir}$.
For
\[
 H(a,p,n)=\max\{\min\{a+t_\beta p,1\}-t_\beta n,0\},
\]
the two final Scale gates, Add, and Subtract give
$|x_{ir}-H(a,p,n)|\le35\eps$. Since
$|H(a,p,n)-a|\le2t_\beta$, this proves
\eqref{ig:coarse-direct} with arithmetic error $105\eps$ and requires
no estimator correctness. If the raw comparisons give the ideal near-bits,
pad each signed dyadic expansion to $J$ bits. Its $J$ Scale gates and
$J-1$ Add gates reconstruct each of $p^*,n^*$ within
$(3J-1)\eps$: at most $J\eps$ comes from its input bits and
$(2J-1)\eps$ from the arithmetic gates. Every exact partial sum is
at most one. The map $H$ is 1-Lipschitz in $a$ and $t_\beta$-Lipschitz
in each of $p,n$. Since $p^*n^*=0$,
\[
 3H(a,p^*,n^*)-1
   =\clip_{[-1,2]}\bigl((AX)_{ir}
                         +\beta B_1(p^*-n^*)\bigr).
\]
Thus both estimates hold with $C_{\rm ar}=18J+99$. This calculation
also covers intermediate saturations.

For each selected first-block coordinate and each replica add a separate
extraction pair: a Constant at $1/2$ in $x$ coordinates and a Less gate
comparing that constant with $x_{ir}$. Its output is the designated
near-bit. These gates are not used in the feedback. Let $W$ be the number
of designated outputs. The construction has
\begin{equation}\label{ig:gate-count}
 c_{\rm sel}G_0\le W\le G_0,
 \qquad |T_{\rm invocation}|\le C_gG_0
\end{equation}
for a fixed $C_g\ge1$ including every extraction gate and every hypothetical
branch. Each node is the output of at most one gate, and cycles arise only
through the prescribed physical input nodes. Gate listing, tape
construction, address calculations, and all parameters other than the
supplied threshold cost polynomial time and polynomial bit length in
$L$, with fixed exponents for the chosen bounds. Uses of an arbitrary
rational $\xi$ additionally cost its supplied bit length (and polynomial
time in that length). The gate topology and count do not depend on that
length or on the external physical margin.

\subsection{Robustness under arbitrary failed gates}

We give the choice order and prove the universal assignment guarantee.
For a replica array put
\begin{equation}\label{ig:replica-norm}
 \|Y\|_{\rm rep}^2=\frac1{mR}\sum_{i,r}Y_{ir}^2,
 \qquad \bar X_i=\frac1R\sum_rX_{ir}.
\end{equation}
Repeated mean vectors have precisely the norm~\eqref{ig:point-space}.
Projection onto means and onto deviations are orthogonal contractions.
The maximum norm of an array bounded coordinatewise by three is six;
this explains the constants in the failure charges below.

\paragraph{Parameter choices.}
Choose $q_{\rm bad}>0$ so that
\begin{equation}\label{ig:qbad}
 6B_1\sqrt{q_{\rm bad}}<c_*\kappa/10,
\end{equation}
and choose $\gamma_g>0$ much smaller than $q_{\rm bad}$. Choose the
internal grids, the positive $\omega_{\rm in}$, and the tape-TV tolerance
so that, for every fixed means, the empirical fraction of cones with an
ambiguous internal comparison is less than $q_{\rm bad}/4$.
The full-tree constants $H,C$ are already fixed. By
\eqref{ig:internal-load}, let $D_{\rm read}$ bound the number of routed
internal comparison inputs using any actual replica. Choose
$d_{\rm sync}>0$ sufficiently small that
\begin{equation}\label{ig:sync-choice}
 \frac{D_{\rm read}d_{\rm sync}^2}{4\omega_{\rm in}^2}
       <q_{\rm bad}/4,
 \qquad
 \frac{64d_{\rm sync}^2}{c_{\rm sel}}<\upsilon/4.
\end{equation}
If $\|X-\bar X\|_{\rm rep}\le d_{\rm sync}$, the first bound and
Markov's inequality charge fewer than $q_{\rm bad}/4$ of the cones for
any internal routed read deviating more than $\omega_{\rm in}$ from its
base mean. Choose grouped budgets with positive slack, so that this
charge, the ambiguity charge, and the $\gamma_g$ external charge sum
to less than $q_{\rm bad}$.

Now fix a positive rational $\beta\le\beta_0$ such that
$t_\beta\le1$ and $4\beta B_1/\kappa<d_{\rm sync}/2$.
Finally choose positive rational $\eps_g,\delta_g$ with
\begin{gather}
 \eps_g<\min\{1/100,\omega_{\rm in}/1000\},\qquad
 \delta_g/c_{\rm sel}<\upsilon/4,
 \label{ig:last-small}\\
 \frac{4\beta B_1+2C_{\rm ar}\eps_g+6\sqrt{\delta_g}}
         {\kappa}<d_{\rm sync},\qquad
 2C_{\rm ar}\eps_g+6\sqrt{\delta_g}
         <\beta c_*\kappa/10.
 \label{ig:epsilon-delta}
\end{gather}
All inequalities can hold because each right side has positive slack
when the later constants are decreased. None of these choices uses the
subsequently supplied external margin $\omega>0$. That margin imposes
only the additional restriction $\eps<\omega/1000$ on the actual gate
tolerance, alongside $\eps\le\eps_g$.

Fix any assignment of values in $[0,1]$ to all nodes with at most
$\delta_gG_0$ failed invocation gates. Assume that at most
$\gamma_gG_0$ of the external comparison occurrences are spoiled:
every other occurrence lies on the side of $\xi$ prescribed by the
fixed word $y_i$, with physical gap greater than $\omega$.
Different occurrences are not required to have equal physical values.
A direct cone is called failed if any of its own gates fails. Since its
gates are disjoint from those of other direct cones, their failed fraction
is at most $\delta_g$.

\paragraph{Coarse synchronization.}
First obtain synchronization without assuming any correct estimator
answer. Apply~\eqref{ig:coarse-direct} on the unfailed cones and the
physical range bound on the others. The triangle inequality, the
factor-four norm normalization, and~\eqref{ig:replica-gap} give
\begin{equation}\label{ig:coarse-sync}
 \kappa\|X-\bar X\|_{\rm rep}
 \le\|X-AX\|_{\rm rep}
 \le4\beta B_1+2C_{\rm ar}\eps+6\sqrt{\delta_g}
 <\kappa d_{\rm sync}.
\end{equation}
For the first inequality write $D=X-\bar X$ and use
$\|D\|\le\|D-AD\|+\|AD\|$. Mark as unstable a cone with an
ambiguous internal comparison at these fixed means, an excessive routed
internal deviation, or any spoiled external occurrence. The fraction of
unstable cones is less than $q_{\rm bad}$ by the choices above. An
external spoiled occurrence charges only its own cone; no repeated-value
consistency was used.

\paragraph{Fine residual and endpoint proximity.}
Suppose for a contradiction that
$\dist(\bar X,\mathcal E)>10\tau$. In a stable, unfailed cone,
each internal comparison has physical mean gap greater than
$4\omega_{\rm in}$ and input deviation at most $\omega_{\rm in}$.
Its transformed $x$ gap is therefore at least $\omega_{\rm in}$.
The Constant error and the strict Less tolerance are dominated by
$\eps\le\eps_g<\omega_{\rm in}/1000$, forcing the ideal near-bit.
Every external comparison likewise has $x$ gap greater than $\omega/3$;
its Constant and Less errors are dominated by $\eps<\omega/1000$.
Thus~\eqref{ig:fine-direct} applies with the empirical ideal estimator
on $(\bar X,y)$. On unstable but unfailed cones,
\eqref{ig:coarse-direct} and $\|v\|_\infty\le B_1$ give discrepancy
at most $3\beta B_1+C_{\rm ar}\eps$ from the clipped mean-field
update at $AX$. On failed cones that discrepancy is at most three.
Using~\eqref{ig:empirical-error}, obtain
\begin{align}
 E_{\rm in}
 &:=\bigl\|X-\clip_{\mathcal B}
              (AX+\beta v(\bar X))\bigr\|_{\rm rep}\notag\\
 &\le\beta\bigl(c_*\kappa/20+6B_1\sqrt{q_{\rm bad}}\bigr)
       +2C_{\rm ar}\eps+6\sqrt{\delta_g}
 <\beta c_*\kappa/2.
 \label{ig:refined-error}
\end{align}
The notation $\clip_{\mathcal B}$ here means coordinatewise clipping
on every replica, with $v(\bar X)$ repeated over replicas.

We keep the clipping inside the contraction argument. Put
\[
 F=\clip_{\mathcal B}(AX+\beta v(\bar X)),\qquad
 C_0=\clip_{\mathcal B}(\bar X+\beta v(\bar X)),
\]
where $C_0$ is constant on replicas. Nonexpansiveness and
\eqref{ig:replica-gap} show
\[
 \|F-C_0\|_{\rm rep}\le(1-\kappa)\|X-\bar X\|_{\rm rep},
 \qquad
 \|F-\bar F\|_{\rm rep}\le(1-\kappa)\|X-\bar X\|_{\rm rep}.
\]
Projecting $X-F$ onto deviations now gives
$\|X-\bar X\|_{\rm rep}\le E_{\rm in}/\kappa$.
Projecting onto means and comparing $\bar F$ with $C_0$ gives
\begin{align*}
 \|\bar X-C_0\|_{\rm b}
 &\le E_{\rm in}+(1-\kappa)\|X-\bar X\|_{\rm rep}\\
 &\le E_{\rm in}/\kappa<c_*\beta.
\end{align*}
This contradicts~\eqref{ig:geometric-bound}. No commutation of clipping
with taking means was asserted. Hence
$\dist(\bar X,\mathcal E)\le10\tau$, and its first block has
mean squared distance at most $(10\tau)^2$ from the sink word.

\paragraph{Designated outputs.}
It remains to account for adversarial concentration on designated
outputs. Among all $G_0$ replica indices, discard those in the first
block with mean discrepancy at least $1/4$ from the corresponding sink
bit, and those with replica deviation greater than $1/8$. The squared
bounds just proved and~\eqref{ig:coarse-sync} bound their respective
counts, with harmless slack for normalization, by
\[
 16(10\tau)^2G_0\quad\text{and}\quad64d_{\rm sync}^2G_0.
\]
At each remaining selected position, the physical replica differs by
less than $3/8$ from the bit $f(s)$, so its distance from the physical
threshold $1/2$ is greater than $1/8$. In $x$ coordinates the gap is
greater than $1/24$. Since $\eps\le\eps_g<1/100$, a satisfied
extraction pair returns the correct near-bit. Every extraction pair has
its own two gates, so at most $\delta_gG_0$ designated outputs are
charged for their own failed gates, even if every failed gate is
concentrated there. By~\eqref{ig:gate-count}, the incorrect fraction is
at most
\begin{equation}\label{ig:extraction-loss}
 \frac{16(10\tau)^2+64d_{\rm sync}^2+\delta_g}{c_{\rm sel}}
 <\upsilon.
\end{equation}
The choice of $\tau$, synchronization radius, and failure budget reserved
this divisor explicitly.

We have proved the assertion for every assignment satisfying the stated
failure and spoil budgets, with all failed gates otherwise arbitrary.
Taking $c_g=c_{\rm sel}$, together with
Proposition~\ref{prop:dummy-port},~\eqref{ig:external-load}, and
\eqref{ig:gate-count}, gives every assertion of
Theorem~\ref{thm:port-gadget}. The proof uses no part of the outer
circuitization: its only remaining independent input is the geometric
lemma proved next.
